% Native semantic reconstruction of source physical page 19 (printed page 18).
% Provenance: IEEE Std 1516-2025, Clause 3.1, definitions designator--federate application.
\HLAStartSemanticPage{19}{18}
\setlength{\HLASemanticDefinitionGap}{13.60pt}
\fontsize{10.02pt}{11.524pt}\selectfont

\begin{hladefinition}[0pt]{def-designator}{designator}
A generic view of arguments referenced in service descriptions. This view improves clarity in situations when arguments (mostly identifiers) to services have different views (implementations) due to a particular programming language or application programmer's interface (API).
\end{hladefinition}

\begin{hladefinition}{def-dimension}{dimension}
A named interval. The interval is defined by an ordered pair of values, the first being the dimension lower bound and the second being the dimension upper bound. A runtime infrastructure (RTI) provides a predefined interval, whose lower and upper bounds are fixed as (0, the dimension's upper bound) as specified in the federation object model (FOM) Document Data (FDD). This interval provides a single basis for communication of all dimension-related data with an RTI. All normalized intervals communicated to the RTI will be subsets of this interval.
\HLADefinitionNote{See 9.1 of IEEE Std 1516.1\HLAStandardMark-2025.}
\end{hladefinition}

\begin{hladefinition}{def-directed-interaction}{directed interaction}
An explicit action taken by a federate directed to a specific object instance that may have some effect or impact on another federate within a federation execution that owns attributes of that object instance.
\end{hladefinition}

\begin{hladefinition}{def-directed-interaction-class}{directed interaction class}
An association between an interaction class and an object class where instances of the interaction class may be directed to be received by an interested federate that owns attribute instances of the object class.
\end{hladefinition}

\begin{hladefinition}{def-discover}{discover}
To receive an invocation of the \HLAService{Discover Object Instance} \HLAPrintersDagger\footnotemark[13] service for a particular object instance.
\HLADefinitionNote{See 6.9 of IEEE Std 1516.1-2025.}
\end{hladefinition}

\begin{hladefinition}{def-discovered-class}{discovered class}
The class that was an object instance's candidate discovery class at a joined federate when that object instance was discovered by that joined federate. \HLASeeAlso{\HLATermReference{candidate discovery class}.}
\HLADefinitionNote{See 5.1.2 in IEEE Std 1516.1-2025.}
\end{hladefinition}

\begin{hladefinition}{def-dynamic-link-compatibility-dlc}{dynamic link compatibility (DLC)}
The ability to interchange the implementation of an API with another without having to recompile or relink the application software that uses the API.
\end{hladefinition}

\begin{hladefinition}{def-epsilon-time-interval}{epsilon time interval}
The distance between any two consecutive logical times T1 and T2 (i.e., the set of logical times between T1 and T2 is empty). For some representations of local time (e.g., floating point), the distance may not be the same for all such instances of T1 and T2.
\end{hladefinition}

\begin{hladefinition}{def-exception}{exception}
Notification of any irregularity that may occur during a service invocation.
\end{hladefinition}

\begin{hladefinition}{def-fault}{fault}
An event that prevents the entire federation from executing in a High Level Architecture (HLA) compliant manner.
\end{hladefinition}

\begin{hladefinition}{def-federate}{federate}
An application that may be or is currently coupled with other software applications under a federation object model (FOM) Document Data (FDD) and a runtime infrastructure (RTI). \HLASeeAlso{\HLATermReference{federate application}\HLASemicolon{} \HLATermReference{joined federate}.}
\end{hladefinition}

\begin{hladefinition}{def-federate-application}{federate application}
An application that supports the High Level Architecture (HLA) interface to a runtime infrastructure (RTI) and that is capable of joining a federation execution. A federate application may join the same federation execution multiple times or may join multiple federation executions. However, each time a federate application joins a federation execution, it creates a new joined federate. \HLASeeAlso{\HLATermReference{joined federate}.}
\end{hladefinition}

\footnotetext[13]{All RTI-initiated services are denoted with a \HLAPrintersDagger{} (printer's dagger) after the service name.}
\clearpage
